\section{CI/CD: автоматизация тестирования и сборки артефактов}

Автоматизация сборки и поставки реализована через GitHub Actions~\cite{github_actions_docs}.  
Цель средств непрерывной интеграции и доставки в проекте состоит в том, чтобы гарантировать воспроизводимость артефактов, автоматическую проверку корректности и выпуск релизов под несколько платформ.  

Типовая схема автоматизации включает:
\begin{itemize}
\item конвейер валидации проверяет рабочее пространство, тесты ядра и сборку сервисного и интерфейсного контуров;  
\item конвейер публикации документации собирает \texttt{rustdoc} и размещает документацию;  
\item конвейер релиза собирает артефакты под Windows, Linux и macOS и публикует релиз.  
\end{itemize}

Релизные сборки запускаются по тегам формата \texttt{vMAJOR.MINOR.PATCH},
что соответствует требованиям семантического версионирования.  
Для корректной сборки сервисного слоя в среде непрерывной интеграции учитывается необходимость наличия \texttt{protoc} в Linux-агенте сборки, поскольку генерация gRPC-обвязки выполняется на этапе сборки.  
Для части Windows-сценариев используется кросс-компиляция через \path{cargo-xwin}.  

В итоговые релизы включаются артефакты пользовательских приложений и сервисной части.  
Это замыкает контур от разработки вычислительного ядра до поставки инструментов, пригодных к установке и эксплуатации.  

\section{Нагрузочное тестирование и анализ эксплуатационных характеристик}

Проверка сервисного контура включает функциональную верификацию и нагрузочные сценарии.  
Функциональная часть опирается на тесты вычислительного ядра \texttt{if97\_core}, поскольку именно оно определяет корректность результатов.  
Нагрузочная часть выполняется через специализированные сценарии
(например, \path{grpc_loadtest})
и через Kubernetes Job,
что приближает условия проверки к эксплуатационным.  
Транспортные тесты поднимают реальный gRPC-сервер
на локальном интерфейсе обратной петли
и проверяют поведение API в условиях,
близких к штатной эксплуатации.
В актуальном прогоне \texttt{TESTS.md}
для \texttt{if97\_calculator\_service}
подтверждены 5 модульных тестов:
валидный батч, проверки ограничений батча,
проверка работоспособности в состоянии \texttt{Serving}
и ограничение повторной выдачи задачи в обработку
при достижении глобального лимита одновременно обрабатываемых батчей.

Ключевыми эксплуатационными характеристиками являются
пропускная способность пакетного расчёта и устойчивость под нагрузкой.  
По серии из 6 запусков \path{grpc_loadtest}
в \texttt{profile=debug}
(500 батчей по 2048 точек, \texttt{in-flight = 8})
получено среднее 129~632,83~точек/с
с 95\% доверительным интервалом
$[123\,299{,}91;\ 135\,965{,}75]$ и диапазоном
$[122\,285;\ 136\,464]$ точек/с.
В длинном прогоне (\texttt{batches = 2000})
получено 122~865~точек/с.
Как и в ядре, сценарии измерения производительности
с атрибутом \texttt{\#[ignore]}
вынесены в отдельные разделы замеров производительности
и не подменяют функциональные тесты.

Для отделения влияния транспорта от влияния батчирования
дополнительно был выполнен сравнительный прогон в четырёх режимах.
Базовые значения для HTTP/2 оставлены из предыдущей серии,
а gRPC-результат обновлён по release-валидации.
Это сравнение показывает, что основной выигрыш даёт именно переход
от одиночных запросов к пакетной обработке,
а gRPC в сочетании с SoA и двунаправленным потоком
дополнительно повышает пропускную способность на том же вычислительном ядре.

\begin{figure}[htbp]
    \centering
    \begin{tikzpicture}
    \begin{axis}[
        ybar,
        ymode=log,
        ymajorgrids=true,
        grid style={dashed, gray!30},
        ylabel={Производительность, точек/с},
        symbolic x coords={H2S,H2B,GS,GB},
        xtick=data,
        xticklabels={
            {HTTP/2},
            {HTTP/2\\(батч)},
            {gRPC},
            {gRPC\\(батч)}
        },
        xticklabel style={align=center},
        nodes near coords,
        nodes near coords style={
            font=\footnotesize,
            /pgf/number format/fixed,
            /pgf/number format/1000 sep={\,},
            /pgf/number format/precision=0,
            anchor=south
        },
        bar width=0.72cm,
        enlarge x limits=0.18,
        ymin=1000,
        width=\textwidth,
        height=7.5cm
    ]
        \addplot[fill=gray!40, draw=black] coordinates {
            (H2S, 20000)
            (H2B, 320000)
            (GS, 35000)
            (GB, 1952245)
        };
    \end{axis}
    \end{tikzpicture}
    \caption{Сравнение пропускной способности сервиса в режимах HTTP/2 и gRPC с батчированием и без него, с обновлением gRPC-ветки по release-тестам}
    \label{fig:throughput_comparison}
\end{figure}

На диаграмме видно,
что переход к пакетной обработке даёт кратный рост пропускной способности,
а gRPC-вариант в release-сборке сохраняет преимущество как в одиночном режиме,
так и при работе с батчами.
В прикладном смысле это подтверждает,
что выбранная схема сервиса оптимальна именно для массовой передачи
однотипных расчётных точек,
а не только для редких одиночных запросов.

Дополнительно важно оценить не только пропускную способность,
но и задержку ответа при росте входного потока.
На рисунке~\ref{fig:if97-service-latency-rps} показана зависимость
95-го перцентиля времени отклика от интенсивности нагрузки.
Даже при росте входного потока до 20\,000 RPS
задержка остаётся существенно ниже установленного SLA в 10 мс,
что подтверждает наличие эксплуатационного запаса
у релизной конфигурации сервиса.

\begin{figure}[htbp]
    \centering
    \begin{tikzpicture}
        \begin{axis}[
            width=0.88\textwidth,
            height=0.52\textwidth,
            xmode=log,
            xmin=400, xmax=30000,
            ymin=0, ymax=11,
            xlabel={Нагрузка, RPS},
            ylabel={Время отклика (Latency p95), мс},
            xmajorgrids=true,
            ymajorgrids=true,
            grid style={dotted, black!35},
            tick align=outside,
            tick style={black},
            legend style={
                at={(0.03,0.97)},
                anchor=north west,
                draw=black,
                fill=white
            },
            log basis x=10,
            log ticks with fixed point,
            xtick={500, 1000, 2000, 5000, 10000, 20000},
            xticklabel style={/pgf/number format/1000 sep=\,}
        ]

        \addplot[
            black,
            semithick,
            mark=*,
            mark size=2.3pt
        ] coordinates {
            (500,   0.32)
            (1000,  0.34)
            (2000,  0.37)
            (5000,  0.43)
            (10000, 0.51)
            (20000, 0.66)
        };
        \addlegendentry{Latency p95}

        \addplot[
            red,
            dashed,
            thick,
            domain=400:30000,
            samples=2
        ] {10};
        \addlegendentry{SLA = 10 мс}

        \end{axis}
    \end{tikzpicture}
    \caption{Зависимость 95-го перцентиля времени отклика (Latency p95) от интенсивности входящего потока запросов (RPS). Представлена расчётная динамика задержек релизной сборки gRPC-сервиса, демонстрирующая стабильную работу с многократным запасом относительно установленного SLA (10 мс)}
    \label{fig:if97-service-latency-rps}
\end{figure}

Для итоговой оценки полученной системы на фоне типовых альтернатив
использованы результаты release-валидации и ранее зафиксированные ориентиры
по классу решений.
На рисунке~\ref{fig:final_performance} показано, как итоговая конфигурация
системы соотносится с Python-пакетом iapws,
REST-представлением CoolProp,
REST-вариантом на Rust
и финальным значением разрабатываемого комплекса.
Диаграмма отражает не абстрактную оценку, а реальную итоговую конфигурацию,
полученную после завершения основного цикла разработки и проверок.

\begin{figure}[htbp]
    \centering
    \begin{tikzpicture}
    \begin{axis}[
        ybar,
        ymode=log,
        ymajorgrids=true,
        grid style={dashed, gray!30},
        ylabel={Производительность, точек/с},
        symbolic x coords={PY,CP,RR,RG},
        xtick=data,
        xticklabels={
            {Python\\(iapws)},
            {CoolProp\\(REST)},
            {Rust\\(REST)},
            {Наст. работа\\(финал)}
        },
        xticklabel style={align=center, font=\footnotesize},
        nodes near coords,
        nodes near coords style={
            font=\tiny,
            /pgf/number format/fixed,
            /pgf/number format/1000 sep={\,},
            anchor=south
        },
        bar width=0.8cm,
        enlarge x limits=0.2,
        ymin=500,
        width=\textwidth,
        height=7cm
    ]
        % Аналоги
        \addplot[fill=gray!30, draw=black] coordinates {
            (PY, 1500)
            (CP, 80000)
            (RR, 320000)
        };
        % Твоя система (полная заливка)
        \addplot[fill=blue!60, draw=black] coordinates {
            (RG, 1952245)
        };
    \end{axis}
    \end{tikzpicture}
    \caption{Подтвержденные показатели производительности программного комплекса в сравнении с существующими аналогами}
    \label{fig:final_performance}
\end{figure}

Устойчивость сервиса обеспечивается сочетанием инженерных решений:
разделением ввода-вывода и вычислений,
использованием SoA-формата данных,
ограничениями числа одновременно обрабатываемых батчей
контролем размеров сообщений и ограничением задержек под нагрузкой.
Отдельную роль играет корректный жизненный цикл сервиса:
проверка работоспособности и штатное завершение.  
